% Standalone problem document, generated from the adjacent README.md.
% Compile with XeLaTeX. Mathematical content is maintained in Markdown.
\documentclass[11pt,a4paper]{article}
\usepackage[fontsize=10.5pt]{scrextend}
\usepackage[margin=22mm,headheight=15pt,headsep=9mm,footskip=11mm]{geometry}
\usepackage{fontspec}
\setmainfont{texgyrepagella}[Extension=.otf,UprightFont=*-regular,BoldFont=*-bold,ItalicFont=*-italic,BoldItalicFont=*-bolditalic]
\setsansfont{texgyreheros}[Extension=.otf,UprightFont=*-regular,BoldFont=*-bold,ItalicFont=*-italic,BoldItalicFont=*-bolditalic]
\setmonofont{texgyrecursor}[Extension=.otf,UprightFont=*-regular,BoldFont=*-bold,ItalicFont=*-italic,BoldItalicFont=*-bolditalic,Scale=MatchLowercase]
\usepackage{amsmath,amssymb}
\usepackage{unicode-math}
\setmathfont{texgyrepagella-math.otf}
\usepackage{xcolor,microtype,enumitem,needspace,fancyhdr}
\usepackage{bookmark}
\definecolor{ink}{HTML}{17344B}
\definecolor{accent}{HTML}{197D80}
\definecolor{muted}{HTML}{596773}
\hypersetup{colorlinks=true,linkcolor=accent,urlcolor=accent,pdftitle={TR-09 - Subquadratic
overparameterization for iterative decomposition of smoothed
tensors},pdfauthor={Open Problems in Numerical Linear Algebra}}
\urlstyle{same}
\setlength{\parindent}{0pt}
\setlength{\parskip}{5pt plus 1pt minus 1pt}
\linespread{1.04}
\setlength{\emergencystretch}{3em}
\widowpenalty=10000
\clubpenalty=10000
\displaywidowpenalty=10000
\allowdisplaybreaks[1]
\setlist{nosep,leftmargin=1.5em}
\providecommand{\tightlist}{\setlength{\itemsep}{0pt}\setlength{\parskip}{0pt}}
\providecommand{\pandocbounded}[1]{#1}
\pagestyle{fancy}
\fancyhf{}
\fancyhead[L]{\sffamily\footnotesize\color{muted}OPEN PROBLEMS IN NUMERICAL LINEAR ALGEBRA}
\fancyhead[R]{\sffamily\bfseries\footnotesize\color{accent}TR-09}
\fancyfoot[L]{\parbox[b]{0.9\textwidth}{\sffamily\fontsize{8}{10}\selectfont\color{muted}Editorial ratings; a bounded search cannot certify that a problem remains unsolved.\\Literature check: 2026-09-08}}
\fancyfoot[R]{\sffamily\footnotesize\color{muted}\thepage}
\renewcommand{\headrulewidth}{0.3pt}
\renewcommand{\headrule}{\hbox to\headwidth{\color{accent}\leaders\hrule height \headrulewidth\hfill}}
\makeatletter
\renewcommand{\section}{\Needspace{5\baselineskip}\@startsection{section}{1}{0pt}{12pt plus 2pt minus 2pt}{4pt}{\sffamily\bfseries\large\color{ink}}}
\renewcommand{\subsection}{\Needspace{4\baselineskip}\@startsection{subsection}{2}{0pt}{9pt plus 1pt minus 1pt}{3pt}{\sffamily\bfseries\normalsize\color{ink}}}
\makeatother
\setcounter{secnumdepth}{0}
\begin{document}
{\sffamily\small\color{muted}Tensor computations\par}
\vspace{3pt}
{\raggedright\hyphenpenalty=10000\exhyphenpenalty=10000\sffamily\bfseries\fontsize{21}{24}\selectfont\color{ink}Subquadratic
overparameterization for iterative decomposition of smoothed
tensors\par}
\vspace{7pt}
{\sffamily\small\textbf{Difficulty:} extreme\quad\textbf{Importance:} broadly
interesting\par}
\vspace{4pt}
{\color{accent}\hrule height 0.8pt}
\vspace{6pt}
\textbf{Status:} open; source published in 2026, with expressly partial
progress in August 2026.

\subsection{Problem statement}\label{problem-statement}

Given \(n\ge r\), take arbitrary base factors
\(\bar A,\bar B,\bar C\in\mathbb R^{n\times r}\), and form \(A,B,C\) by
adding independent Gaussian entries of variance \(\rho^2/n\), where
\(\rho=r^{-q}\) and \(q>0\) is fixed. The input is

\[
T=\sum_{j=1}^{r}a_j\otimes b_j\otimes c_j.
\]

Does there exist \(k(r)=o(r^2)\), a polynomial lower bound \(n_0(r)\),
and an explicitly specified ALS or gradient-descent method, initialized
randomly, such that for every \(n\ge n_0(r)\) and \(0<\varepsilon<1\) it
returns \(k\) summands with

\[
\left\|T-\sum_{i=1}^{k}x_i\otimes y_i\otimes z_i\right\|_F
\le\varepsilon\|T\|_F
\]

using polynomially many exact real arithmetic operations in
\(n,r,\log(1/\varepsilon)\), with probability \(1-o_r(1)\) over the
once-drawn smoothed input and at least inverse-polynomial conditional
success over initialization? Update rules, initialization, and permitted
restarts must be specified. Initial factors must be sampled
independently of the input and its unknown factors, apart from an
explicitly specified scalar scale normalization and dependence on
\(n,r,\rho,\varepsilon\). The method must optimize the squared residual
by alternating block least squares or gradient steps; substituting an
algebraic tensor-recovery algorithm does not answer the question. The
base-factor quantifier must be preserved rather than replaced by
assumptions about the successful trajectory.

\subsection{References}\label{references}

Dionysis Arvanitakis, Vaidehi Srinivas, and Aravindan Vijayaraghavan,
\href{https://proceedings.mlr.press/v336/arvanitakis26a.html}{\emph{Open
Problem: How Much Overparametrization Is Needed for ALS in Tensor
Decomposition?}}, COLT 2026, PMLR 336, 7105--7110, §2, Open Problem 2
and the following probability clarification. Zhang et al.,
\href{https://arxiv.org/html/2608.13060v1}{\emph{VALG: An Agentic System
for ML Theory Research and Demonstrations on COLT 2026 Open Problems}},
§4.1.1, Theorem 4.1 and Discussion.

\subsection{Status check}\label{status-check}

Searches included \textit{"ALS" "subquadratic" tensor 2026} and the
exact COLT problem title. VALG reports
\(k=\Theta(r^{5/3}(\log r)^{5/2})\) under additional scale,
interference, balance, smoothing, and dimension assumptions, and
explicitly calls the result partial. That result does not resolve the
general base-factor formulation. The source informally motivates
well-conditioned factors but its displayed smoothed model permits
arbitrary base factors; this entry follows that displayed quantifier. A
future revision may distinguish a uniformly norm-bounded base model once
a precise source formulation is identified.
\end{document}
